% ============================================================
%  sections/conclusions.tex  –  Conclusions
% ============================================================

This paper provides the first transaction-level estimate of the price elasticity of marijuana demand in Mexico.  Using 3{,}293 crowdsourced purchases observed across all 32 states and 351 municipalities, the preferred state-by-month fixed-effects specification yields an elasticity of $-0.652$.  The estimate implies that a 10\% higher unit price is associated with a 6.5\% smaller quantity purchased in a completed transaction.  This is an intensive-margin estimate: it speaks to transaction size among observed buyers, not to participation, purchase frequency, total individual consumption, or aggregate market demand.  Within that margin, the result is economically meaningful but less than unit-elastic, close to transaction-based evidence from other cannabis markets \citep{davis2016crowd,halcoussis2017geo,riley2020southafrica} and within the broader range summarized by \citet{gallet2014monkey}.

A central finding is that product quality matters.  Seedless marijuana is substantially more expensive and less price responsive, with an estimated elasticity of about $-0.61$, while seeded marijuana is more elastic, at about $-0.89$.  This difference suggests that consumers do not treat all marijuana transactions as a homogeneous product.  For policy design, the quality gradient matters because weight-based, price-based, and potency-sensitive tax instruments may affect product tiers differently.  A uniform per-gram tax, for example, would represent a larger proportional price increase for lower-quality product, precisely the segment in which demand appears more price responsive.

The benchmark estimate is stable across the main sensitivity checks.  Trimming extreme price and quantity observations changes the elasticity only modestly, excluding visibly heaped quantities preserves a negative and significant estimate, and adding municipality fixed effects yields an elasticity of $-0.664$ in the restricted sample.  Wild cluster bootstrap inference also confirms that the main result is not an artifact of having only 32 state clusters.  The Oster coefficient-stability exercise reports $|\delta^{*}|=7.70$; because the raw delta is negative under the sign convention induced by the negative coefficients, this magnitude should be read as the strength of unobserved selection required to move the benchmark estimate to zero.  These exercises do not prove causal identification, but they make a purely spurious interpretation of the price--quantity relationship less plausible.

The instrumental-variable evidence should be read in the same spirit.  The OSM production-hub driving-time instrument remains too weak to be informative as a standalone source of identification.  The hub-only one-month-lag marijuana-seizure gravity instrument is the closest IV analogue to the benchmark: it clears the conventional first-stage threshold and yields an elasticity of $-0.653$, while a stricter outside-region seizure instrument weakens substantially.  External Hausman instruments that exclude the buyer's own state or macroregion have strong first stages and produce negative elasticities of similar order, but they still rely on the assumption that other-market price variation captures supply-side cost shocks rather than correlated demand shocks.  The IV block therefore triangulates the benchmark and helps bound plausible interpretations, but it is not presented as a definitive causal proof.

The policy implication is not that a particular Mexican cannabis tax rate follows mechanically from an elasticity of $-0.65$.  The estimate comes from illicit-market transactions before national recreational regulation, not from a legal market with observed pass-through, compliance costs, product standards, or substitution between legal and illegal suppliers.  Still, it is a necessary empirical input.  Less-than-unit elastic demand implies that price-based regulation can reduce transaction quantities, but the response is unlikely to be so large that the tax base disappears.  At the same time, the quality-tier heterogeneity cautions against designing cannabis taxes as if the market were a single homogeneous product.  Any regulatory framework should therefore combine elasticity evidence with explicit assumptions about pass-through, enforcement, public-health objectives, and the coexistence of legal and illegal supply.

Several limitations remain.  The survey is crowdsourced and was promoted online, so the sample overrepresents young, male, and more educated consumers relative to the full population of marijuana users.  Quantities and prices are self-reported, and unit prices are constructed from reported value and quantity, so measurement error and nonlinear package pricing cannot be ruled out even after trimming and heaping checks.  Future work could combine transaction-level purchase data with representative consumption surveys, regulated-market administrative records if legalization proceeds, or repeated observations of buyers before and after policy changes.  Until such data exist, the evidence here offers a transparent and empirically grounded estimate of marijuana demand in Mexico at the transaction margin most directly relevant for price, taxation, and market-design debates.
